\documentclass[10pt]{article}
\usepackage{amsmath,amssymb,mathtools,mathrsfs}
\usepackage[margin=1in]{geometry}
\emergencystretch=12pt
\newcommand{\im}{\operatorname{im}}
\newcommand{\rank}{\operatorname{rank}}
\begin{document}

% Complete original provider: ORIGINAL_PRIMITIVE_SOURCE_METRIC.tex
% Original primitive source: exact polynomial dual and one-step native transport.
\section{The original primitive source and its adjacent Gamma metric}
\label{sec:PDS}

This calculation continues the complete LCA15--21 polynomial quotient
construction. Its new conclusions are the explicit codimension-three
truncation, its exact rank-one native metric change, and a finite
coefficient-product bound for that change. All spaces and measures below
are the original ones. The connecting target metric and its later
long return remain the literal target metrics; the source computation
here is at the two specified low cutoffs.

\subsection{Original objects and the polynomial quotient map}

Fix the actual simple quartet, its invertible period matrix $\mathsf H$,
logarithm branch, and full conductor coefficients. Retain
\[
\begin{gathered}
 \beta_*={\tfrac12}-\delta-i\gamma,\qquad
 \beta_{n,a,b}=n\beta_*+2a\delta+2bi\gamma,\quad\delta,\gamma>0,\\
 E_A(z)=\sum_{a,b=0}^8a_{ab}e^{\beta_{8,a,b}z},\quad
 \mu_r=E_A^{(r)}(0),\quad
 v=\operatorname{ord}_0E_A\le80,\quad\mu_v\ne0,\\
 L_i(Z,W)=(\mathsf H(ZW,Z,W,1)^{\mathsf T})_i,\quad
 b_{i,a,b}=[Z^aW^b]L_i^4,\quad
 B_i(z)=\sum_{a,b=0}^4b_{i,a,b}e^{\beta_{4,a,b}z},\\
 b_i=B_i(0)=\sum_{a,b=0}^4b_{i,a,b},\qquad
 \sum_{i=1}^4D_i(z)B_i(z)=1 .
\end{gathered}\tag{PDS1}
\]
The $D_i$ are the proved original entire Bezout expressions. In
particular at least one $b_i$ is nonzero. Keep the full biform conductor,
including repeated components and zero edge coefficients. Set
\[
\begin{gathered}
 k=4l+1\ge9,\quad j=k+4,\quad q_k=(k+1)^2,\quad Q=(k+5)^2,\\
 e\in\{0,1\},\quad D_e=Q-v-1+e,\quad
 \ell_e=8k+23+e,\quad d_e=\ell_e+1,\\
 \chi_n(S)=\prod_{a,b=0}^n(S-\beta_{n,a,b}),\qquad
 (\mathcal T_Ap)(S')=\sum_{a,b=0}^8a_{ab}p(S'+\beta_{8,a,b}),\\
 (\mathcal T_{B_i}p)(S)=
       \sum_{a,b=0}^4b_{i,a,b}p(S+\beta_{4,a,b}).
\end{gathered}\tag{PDS2}
\]
All dual functionals are complex-linear. For a finite germ $F$ through
degree $D_e$, let $\Lambda_F(S'^r)=F^{(r)}(0)$.
Let $\mathcal H_e=\mathbb C[S']_{\le D_e}^{\vee}$, and let
$\mathcal S_e\subset\mathcal H_e$ be the jets of the original analytic
germs
\[
 F(z)=\frac{N(z)}{E_A(z)},\qquad
 N(z)=\sum_{a,b=0}^k n_{ab}e^{\beta_{k,a,b}z},\qquad
 N^{(r)}(0)=0\quad(0\le r<v).
\tag{PDS3}
\]
Define the unchanged relation polynomials and quotient coordinates by
\[
 r_a(S')=\mathcal T_A(\chi_k(S)S^a),\quad0\le a\le\ell_e,
 \qquad \pi_e(F)=(\Lambda_F(r_a))_{a=0}^{\ell_e}.
\tag{PDS4}
\]
Then $\pi_e$ is onto and $\ker\pi_e=\mathcal S_e$.
Here is the complete kernel argument. The distinct upper $k$ nodes give
a Vandermonde matrix, whose first $v$ rows have rank $v$ since
$v\le80<q_k$. The numerator space in PDS3 therefore has dimension
$q_k-v$. Its jet map is injective: if $F$ vanishes through degree
$D_e$, then $N=E_AF$ vanishes through degree $D_e+v\ge q_k-1$;
the first $q_k$ Vandermonde rows force every $n_{ab}=0$.
The Leibniz identity gives
$\Lambda_F(\mathcal T_Ap)=\Lambda_{E_AF}(p)$, so all $r_a$ annihilate
PDS3. Finally
\[
 \mathcal T_AS^m=\sum_{r=0}^m\binom mr\mu_r S'^{m-r},\qquad
 \deg r_a=q_k+a-v,\qquad
 \operatorname{lc}(r_a)=\mu_v\binom{q_k+a}{v}.
\tag{PDS5}
\]
These successive degrees prove independence. Their common kernel has
dimension $D_e+1-d_e=q_k-v$, proving the assertion. This is precisely
the LCA16 quotient, with no new row norm. Every original lower
degree-$(k-8)$ exponential is in $\mathcal S_e$, since its product
with $E_A$ has an allowed degree-$k$ numerator. Thus taking the
original lower-gauge quotient first and then quotienting by its
analytic-source image gives exactly the same affine fibres as
$\mathcal H_e/\mathcal S_e$.

\subsection{The full original Gamma dual metric}

The source member of the two original order pairs is $s=1$ or $s=k$,
and its centre is $c=k/2-4$. Keep its entire measure and recurrence:
\[
\begin{gathered}
 M_s=(2\pi)^{s/2},\qquad
 dm_s(y)=\frac{M_s2^{s/2}}{4\pi\Gamma(s/2)}
           |\Gamma(s/4+iy/2)|^2\,dy,\qquad
 h_{n,s}=M_s n!(s/2)_n,\\
 p_{0,s}=1,\quad p_{1,s}=y,\quad
 p_{n+1,s}=yp_{n,s}-n(n-1+s/2)p_{n-1,s},\qquad
 \phi_n(S')=p_{n,s}((S'-c)/i)/\sqrt{h_{n,s}}.
\end{gathered}\tag{PDS6}
\]
These are the original orthonormal polynomials for the norm
$\|p\|_{s,c}^2=\int_{\mathbb R}|p(c+iy)|^2dm_s(y)$,
as established in the original Gamma construction reproduced in the
cumulative providers. The coefficient isometry uses this complete
measure, rather than dividing its mass out. Let $\mathsf B_e$ be
the triangular matrix defined by
$S'^r=\sum_{n=0}^{D_e}(\mathsf B_e)_{nr}\phi_n(S')$ and let
\[
 C_e=\bigl(\mathsf B_e\operatorname{coeff}_{S'}r_a\bigr)_{a=0}^{\ell_e},
 \qquad H_e=C_e^*C_e>0,\qquad G_e=\overline{H_e^{-1}}.
\tag{PDS7}
\]
The columns of $C_e$ are polynomial coefficients in the original
orthonormal frame. For $w=(\Lambda_F(\phi_n))_n$ the literal
quotient row map is $x=C_e^{\mathsf T}w$. Its minimum-norm lift is
$w=\overline C_e\,\overline H_e^{-1}x$ and has squared norm
$x^*G_ex$. Indeed $C_e^{\mathsf T}\overline C_e=\overline H_e$,
and the displayed lift is orthogonal to $\ker C_e^{\mathsf T}$.
This proves both the full complex-linear dual orientation and the
quotient metric, including every conjugation in PDS7. The same metric
results from the two successive minima over the original lower
gauges and analytic source, by the equality of fibres proved after
PDS5.

The degree-one primitive condition at the first endpoint is empty.
At the second endpoint its four original rows are
\[
 R_i(S)=\frac{\mathcal T_{B_i}\chi_j(S)}{\chi_k(S)},\qquad
 W=\bigl(\operatorname{coeff}_{S}R_i\bigr)_{i=1}^4
       \in\mathbb C^{d_1\times4},\qquad
 \mathcal P_0=\mathbb C^{d_0},\quad
 \mathcal P_1=\ker W^{\mathsf T}\subset\mathbb C^{d_1}.
\tag{PDS8}
\]
Here $\mathcal P_e$ denotes the primitive quotient source, not a
biform section space. All degree-$k$ roots become degree-$j$ roots
after every degree-four shift; hence the numerator is divisible by
$\chi_k$ and $\deg R_i\le\ell_1$. Translation sums commute, so
$\mathcal T_A(\chi_kR_i)=\mathcal T_{B_i}\mathcal T_A\chi_j$.
This is exactly the original positive degree-zero exterior
differential, and its row is $W^{\mathsf T}$, as in LCA18--19.
The full independence proof LCA20--21 applies to the actual
invertible period, without a genericity exclusion: the four
$L_i^4$ are independent on the rank-four Segre quadric, and the
25 distinct translations of the degree-$Q-v\ge116$ polynomial
$\mathcal T_A\chi_j$ give an injective Vandermonde moment map.
Thus $\operatorname{rank}W=4$. The complete LCA20--21 proof is
retained as a provider in the cumulative source, rather than
replacing this fact by a new period assumption.

The metric on $\mathcal P_e$ is the restriction of $G_e$.
At $e=1$ its exact polynomial dual is
\[
 \operatorname{span}(r_0,\ldots,r_{\ell_1})\big/
 \operatorname{span}(\mathcal T_A(\chi_kR_i))_{i=1}^4.
\]
To verify the metric as well as the vector-space statement,
choose any fixed full basis matrix $Z_1$ of $\ker W^{\mathsf T}$.
A coefficient vector $f$ of a relation polynomial restricts to the
functional $Z_1^{\mathsf T}f$ on this basis. Minimizing $f^*H_1f$
over all lifts with $Z_1^{\mathsf T}f=a$ gives
\[
 \min_{Z_1^{\mathsf T}f=a}f^*H_1f
 =a^*(Z_1^{\mathsf T}H_1^{-1}\overline Z_1)^{-1}a
 =a^*\overline{(Z_1^*G_1Z_1)^{-1}}a .
\tag{PDS9}
\]
The minimizer follows by the adjoint of
$Z_1^{\mathsf T}H_1^{-1/2}$; subtracting it leaves an orthogonal
kernel vector. The kernel of $Z_1^{\mathsf T}$ is exactly
$\operatorname{im}W$, by dimensions and $W^{\mathsf T}Z_1=0$.
This proves the stated original quotient metric and its exact dual
morphism, with no deletion of the four row polynomials.

\subsection{The exact codimension-three truncation}

The coefficient of $S^{\ell_1}$ in $R_i$ is $b_i$, since both
$\chi_j$ and $\chi_k$ are monic. Choose a fixed original index
$i_*$ with $b_{i_*}\ne0$, write $b_*=b_{i_*}$, and retain
\[
 R_{i_*}(S)=\sum_{a=0}^{\ell_0}(u_*)_aS^a+b_*S^{\ell_1},\qquad
 C_i(S)=R_i(S)-\frac{b_i}{b_*}R_{i_*}(S)\quad(i\ne i_*).
\tag{PDS10}
\]
These displayed combinations are exact coordinates for the
intersection of the original four-row span with lower degree;
the original rows and coefficients in PDS8 remain fixed.
The three $C_i$ have degree at most $\ell_0$ and are independent:
a relation among them would give a relation among all four $R_i$,
whose coefficients on indices other than $i_*$ are unchanged.
Let $V$ have the three coefficient columns $C_i$, in the original
index order with $i_*$ omitted. Then
\[
 \begin{gathered}
 \mathcal Z=\ker V^{\mathsf T}\subset\mathbb C^{d_0},\qquad
 \dim\mathcal Z=d_0-3=8k+21,\qquad
 E=\begin{bmatrix}I_{d_0}\\ -u_*^{\mathsf T}/b_*\end{bmatrix},\\
 \operatorname{pr}_{0}:\mathcal P_1\xrightarrow{\ \simeq\ }\mathcal Z,
 \quad\operatorname{pr}_{0}^{-1}=E|_{\mathcal Z}.
 \end{gathered}
\tag{PDS11}
\]
In fact the $i_*$ row of $W^{\mathsf T}(x_0,x_\mathrm{top})=0$
gives $x_\mathrm{top}=-u_*^{\mathsf T}x_0/b_*$. Substitution in
the other three rows gives precisely $V^{\mathsf T}x_0=0$.
The first $d_0$ relation polynomials in PDS4 are unchanged at the
two cutoffs, so $\operatorname{pr}_0$ is also the actual jet
truncation on the primitive quotient source. If an upper analytic
source is added its truncated jet is an old analytic source; hence
this map is well-defined on the original quotient. Surjectivity
and the displayed inverse follow from the equations just proved.

\subsection{One rank-one update, with every frame retained}

Choose an arbitrary fixed full column basis $Z$ of $\mathcal Z$;
use $EZ$ at the upper endpoint. For the same source order $s$ and
centre $c$, the old orthonormal polynomials are unchanged when the
cutoff increases. In particular the old $C_0$ is padded by a zero
last row in $C_1$. Write the exact block matrix and Schur scalar
\[
 H_1=\begin{bmatrix}H_0&b\\b^*&a\end{bmatrix},\qquad
 b=C_0^*c_\mathrm{top},\quad a=c_\mathrm{top}^*c_\mathrm{top},\qquad
 \sigma=a-b^*H_0^{-1}b>0,
\tag{PDS12}
\]
where $c_\mathrm{top}$ is the full coefficient column of
$r_{\ell_1}$. Padded columns are understood in the products.
Define
\[
 \eta=-u_*^{\mathsf T}/b_*-b^{\mathsf T}\overline H_0^{-1},\quad
 K_0=Z^*\overline H_0^{-1}Z,\quad
 K_1=(EZ)^*\overline H_1^{-1}(EZ).
\tag{PDS13}
\]
Then the complete metric comparison is
\[
 \boxed{\quad K_1=K_0+\frac{(\eta Z)^*(\eta Z)}{\sigma},\qquad
 \frac{\det K_1}{\det K_0}
   =1+\frac{(\eta Z)K_0^{-1}(\eta Z)^*}{\sigma}\ge1.\quad}
\tag{PDS14}
\]
For proof, block inversion of $\overline H_1$ gives for every
$(x_0,x_\mathrm{top})$
\[
 \begin{bmatrix}x_0\\x_\mathrm{top}\end{bmatrix}^{\!*}
 \overline H_1^{-1}
 \begin{bmatrix}x_0\\x_\mathrm{top}\end{bmatrix}
 =x_0^*\overline H_0^{-1}x_0
  +\frac{|x_\mathrm{top}-b^{\mathsf T}
                    \overline H_0^{-1}x_0|^2}{\sigma}.
\tag{PDS15}
\]
Insert PDS11 and then $x_0=Zz$. This proves the matrix equality.
The determinant identity follows by taking the determinant of
$I+K_0^{-1/2}(\eta Z)^*(\eta Z)K_0^{-1/2}/\sigma$;
this rank-at-most-one operator has its only possible nonzero
eigenvalue $(\eta Z)K_0^{-1}(\eta Z)^*/\sigma$.
The frame $Z$ has not been changed: replacing it by $ZT$ changes
both determinants by the same $|\det T|^2$. PDS14 also proves
that the actual truncation $\mathcal P_1\to\mathcal Z$ is a
contraction for the displayed native metrics, and locates its
entire metric defect in this one explicit row.

\subsection{An explicit coefficient-product upper estimate}

Set $D=D_1=Q-v$ and keep the literal original polynomial
\[
 r_*(S')=\mathcal T_A(\chi_kR_{i_*})(S')
         =\mathcal T_{B_{i_*}}\mathcal T_A\chi_j(S'),\qquad
 \lambda_*=\operatorname{lc}r_*=b_*\mu_v\binom Qv\ne0.
\tag{PDS16}
\]
There is an exact finite upper bound
\[
 1\le\frac{\det K_1}{\det K_0}
 \le\frac{\|r_*\|_{s,c}^2}{|b_*|^2\sigma}
 \le\frac{\|r_*\|_{s,c}^2}{|\lambda_*|^2h_{D,s}}.
\tag{PDS17}
\]
To prove the first upper inequality, the orthogonal projector onto
$\operatorname{im}(\overline H_0^{-1/2}Z)$ gives
$ZK_0^{-1}Z^*\le\overline H_0$.
Thus the numerator in PDS14 is at most
$\eta\overline H_0\eta^*$.
Let $v_*=u_*/b_*+H_0^{-1}b$. Then $\eta=-v_*^{\mathsf T}$,
so this upper bound is $v_*^*H_0v_*$. The orthogonal
projection of $r_*$ onto the old relation span has coefficient
vector $b_*v_*$, whereas its orthogonal remainder is $b_*$
times the remainder of $r_{\ell_1}$. Consequently
\[
 \|r_*\|_{s,c}^2=|b_*|^2(v_*^*H_0v_*+\sigma).
\tag{PDS18}
\]
This proves the first inequality of PDS17. Every old relation
has degree at most $D-1$. The coefficient of $\phi_D$ in
$r_{\ell_1}$ is
$i^D\mu_v\binom Qv\sqrt{h_{D,s}}$; this uses its actual
leading coefficient and the phase $i^{-D}/\sqrt{h_{D,s}}$
in PDS6. This coefficient remains in its orthogonal remainder.
Hence $\sigma\ge|\mu_v\binom Qv|^2h_{D,s}$, proving the
second inequality of PDS17 with its original scalar intact.

For a completely finite coefficient bound define
\[
\begin{gathered}
 L=\max\{1,|c|+2(D+1+s/2)\},\qquad
 \Theta_{a,b,r,t;u,w}
       =\beta_{8,a,b}+\beta_{4,r,t}-\beta_{j,u,w},\\
 P_*=
 \sum_{a,b=0}^8\sum_{r,t=0}^4|a_{ab}b_{i_*,r,t}|
          \prod_{u,w=0}^j(L+|\Theta_{a,b,r,t;u,w}|),\qquad
 U_*=
 \frac{M_s P_*^2}
 {|b_*\mu_v\binom Qv|^2h_{D,s}} .
\end{gathered}\tag{PDS19}
\]
The $L$ in this subsection is a numerical bound and is not the
jet length denoted $L$ in LCA2. Then
\[
 0\le\log\frac{\det K_1}{\det K_0}\le\log U_*.
\tag{PDS20}
\]
Here is a direct proof that PDS19 bounds the original Gamma norm.
In the orthonormal polynomial basis in $y$, multiplication by $y$
has raising coefficient $\sqrt{(n+1)(n+s/2)}$ and lowering
coefficient $\sqrt{n(n-1+s/2)}$, by PDS6 and
$h_{n+1,s}/h_{n,s}=(n+1)(n+s/2)$.
On polynomials of degree at most $D-1$ both weighted shifts
have norm at most $D+1+s/2$; their sum has norm at most
$2(D+1+s/2)$. Thus multiplication by $S'=c+iy$, through
degree $D$, has norm at most $L$. Starting with the constant
polynomial, whose squared norm is $M_s$, proves
\[
 \|p\|_{s,c}\le\sqrt{M_s}\sum_{m=0}^D|[S'^m]p|L^m
       \quad (\deg p\le D).
\tag{PDS21}
\]
The exact expansion of PDS16 is
\[
 r_*(S')=\sum_{a,b=0}^8\sum_{r,t=0}^4a_{ab}b_{i_*,r,t}
                  \prod_{u,w=0}^j(S'+\Theta_{a,b,r,t;u,w}).
\tag{PDS22}
\]
The weighted coefficient $\ell^1$ norm is submultiplicative:
for two polynomials its coefficient convolution and the triangle
inequality give $\sum_m|(fg)_m|L^m\le
(\sum_m|f_m|L^m)(\sum_m|g_m|L^m)$.
Applying this identity to every factor in PDS22 bounds the
weighted coefficient sum of $r_*$ by $P_*$. Its coefficients
above degree $D=Q-v$ vanish by the exact moment identities in
PDS5; bounding the full degree-$Q$ expansion before those
cancellations is still an upper bound for the unchanged
degree-$D$ polynomial. PDS21 now gives
$\|r_*\|_{s,c}^2\le M_sP_*^2$, and PDS17 proves PDS20.
In particular $U_*\ge1$ is a consequence, rather than an
extra assumption on the original coefficients.

There is also an explicit uniform growth bound in either
original source-order family. Define the fixed positive numbers
\[
 \begin{gathered}
 R_* = \sqrt{(1/2+\delta)^2+\gamma^2},\quad
 A_{\rm abs}=\sum_{a,b=0}^8|a_{ab}|,\quad
 B_{\rm abs}=\sum_{r,t=0}^4|b_{i_*,r,t}|,\\
 C_*=6+17R_*,\quad
 C_{\rm abs}=\max\left\{1,\frac{2(A_{\rm abs}B_{\rm abs})^2}{|b_*\mu_v|^2}\right\}.
 \end{gathered}
\tag{PDS23}
\]
The original nodes satisfy
$|\operatorname{Re}\beta_{n,a,b}|\le n(1/2+\delta)$ and
$|\operatorname{Im}\beta_{n,a,b}|\le n\gamma$. Both bounds are
attained together at $(a,b)=(n,0)$ and $(n,n)$; thus
$\max_{a,b}|\beta_{n,a,b}|=nR_*$,
and therefore $|\Theta|\le(k+16)R_*$.
For $s=1$ or $s=k$, $D\le Q$, $k\le Q$, and
$|c|\le k/2+4$, so
$L\le 2Q+3k/2+6\le6(Q+1)$.
Also $k+16\le17(Q+1)$. Every factor in PDS19 is consequently
at most $C_*(Q+1)$, giving
$P_*\le A_{\rm abs}B_{\rm abs}[C_*(Q+1)]^Q$.
The full factorial denominator gives a sharper bound than
dropping its growth. For $s\ge1$,
\[
 D!(s/2)_D\ge D!(1/2)_D=\frac{(2D)!}{4^D}
                         \ge\left(\frac{D}{\exp(1)}\right)^{2D}.
\]
The product identity follows by separating the even and odd
factors of $(2D)!$. The last inequality follows from
$\sum_{r=1}^{2D}\log r\ge\int_1^{2D}\log x\,dx
=2D\log(2D)-2D+1$.
Since $D=Q-v$, $Q\ge196$ and $v\le80$ imply
$Q+1\le2D$, the exact PDS19 bound therefore satisfies
\[
 \begin{split}
 U_*&\le
 \frac{(A_{\rm abs}B_{\rm abs})^2}{|b_*\mu_v|^2}
 C_*^{2Q}(Q+1)^{2Q}
                \left(\frac{\exp(1)}D\right)^{2D}\\
 &\le C_{\rm abs}(Q+1)^{2v}
                         (2\exp(1)C_*)^{2Q}.
 \end{split}
\]
Here $\binom Qv\ge1$ was used only in an upper bound; the
full binomial, mass and factorial denominator remain in PDS19.
For the second inequality split $(Q+1)^{2Q}$ into its
$2v$ and $2D$ powers and use $Q+1\le2D$, $D\le Q$ and
$C_*\ge6$. Consequently the proved finite bound is
\[
 \boxed{\quad
 0\le\log\frac{\det K_1}{\det K_0}
 \le\log U_*
 \le\log C_{\rm abs}+2v\log(Q+1)
       +2Q\log(2\exp(1)C_*)
 =O_{\mathrm{actual}}(Q)=o(kq_k).
 \quad}\tag{PDS24}
\]
The last limit uses $Q/q_k\to1$, fixed $v$, and $1/k\to0$ for
$k\equiv1\pmod4$. The constants are fixed by the actual
quartet and period, independent of $k$ and of the choice between
the two original source-order families. The sharp finite PDS19
retains the full mass, phase-dependent nodes, factorials, conductor
coefficients and leading scalar; PDS24 proves the scale of its
upper bound without assuming an inverse estimate.

The maps PDS11 and PDS14 now give the primitive part of the
low-cutoff source comparison inside the original LCA connecting
sequence. The full cohomology additionally has the retained
$\ker U_e$ constituent from LCC. Passing these source metrics
through the actual connecting map requires the original target
quotient Gram; PDS9 specifies the exact source dual to be used
in that calculation, and PDS20--24 control its adjacent source
correction. None of these equalities identifies a later target
coefficient flag with a common compatible source at all cutoffs.


% Complete original provider: ORIGINAL_CONNECTING_IMAGE_TARGET_GRAM.tex
% Complete native target realization of the actual low connecting image.
\section{The original connecting image and its full native target determinant}
\label{sec:CIT}
Keep the complete CNC, LCC and LCA data, including the nonreduced
$A^{\rm hom}$, every original coefficient and the exact
$E_A(z)=\sum a_{ab}e^{\beta_{8,a,b}z}$. Put
\[
 \begin{gathered}
 k=4l+1\ge9,\quad j=k+4,\quad q_a=(a+1)^2,\quad q'_a=(a-7)^2,\\
 e\in\{0,1\},\quad D_e=q_j-v-1+e,\quad L_e=D_e+1,\quad
 0\le v\le80,\quad\mu_v=E_A^{(v)}(0)\ne0,\\
 (s_0,s_1)=(1,1)\ \hbox{or}\ (k,j),\qquad
 c'_a=a/2-4,\qquad m_a=q_a-q'_a-v .
 \end{gathered}\tag{CIT1}
\]
The two source orders in this display remain paired throughout;
write $s_k=s_0$ and $s_j=s_1$ when a degree subscript is used.
All finite functionals are complex-linear. At a fixed $D_e$ let
$\mathcal H_0,\mathcal H_1$ be their original Gamma polynomial
duals in degrees $k,j$, with full orthonormal frames
\[
 \phi_{r,s,c'}(S')=
 \frac{p_{r,s}((S'-c')/i)}
      {\sqrt{(2\pi)^{s/2}r!(s/2)_r}},\qquad 0\le r\le D_e .
 \tag{CIT2}
\]
Use the exact CNC10 matrices $T_0=T_{k,D_e}$, $T_1=T_{j,D_e}$
for its original analytic numerators divided by $E_A$, and
$V_0,V_1$ for every original lower evaluation. Thus
$[V_0,T_0]$ has its original full column order.

\subsection{The literal primitive domain and target coefficient matrix}

The four multiplier matrices are completely specified without an
unspecified evaluator. Let $\mathcal B_a$ be the original
monomial-to-Gamma matrix defined by
$S'^r=\sum_n(\mathcal B_a)_{nr}\phi_{n,s_a,c'_a}$ and set
\[
 \begin{gathered}
 U_a=\operatorname{diag}(1/r!)_{r=0}^{D_e}\mathcal B_a^{\mathsf T},
 \qquad
 (\operatorname{Mult}_{L_e}(B_i))_{rt}=
 \begin{cases}B_i^{(r-t)}(0)/(r-t)!,&r\ge t,\\0,&r<t,\end{cases}\\
 \mathsf T_i=U_j^{-1}\operatorname{Mult}_{L_e}(B_i)U_k,\qquad
 \mathsf T=\operatorname{stack}_{i=1}^4\mathsf T_i,\qquad
 \mathcal V=\operatorname{diag}(V_1,4),\quad
 \mathcal T_1=\operatorname{diag}(T_1,4).
 \end{gathered}\tag{CIT3}
\]
Here $\operatorname{diag}(P,4)$ means four copies. The map $U_a$
sends the actual native functional values to Taylor coefficients:
ordinary derivatives are $\mathcal B_a^{\mathsf T}w$.
This proves CIT3 and preserves every factorial and the original
leading phase $i^r$ inside $\mathcal B_a$. Its full inverse tuple
is the analogous multiplication by the original entire $D_i$;
$\sum D_iB_i=1$ gives $\mathsf S\mathsf T=I$.

Define the raw source analytic space
\[
 \mathcal S^{\rm raw}_k=\operatorname{im}[V_0,T_0]\subset\mathcal H_0,
 \qquad \mathcal L_0=\operatorname{im}V_0,\quad
 \mathcal L_1=\operatorname{im}\mathcal V .
 \tag{CIT4}
\]
It is exactly the set of native functionals of $N/E_A$ for all
original upper numerators $N$ vanishing to order $v$. Indeed
$T_0$ gives the original DNE quotient frame and $V_0$ gives exactly
its allowed numerator changes by conductor multiples. The columns
are independent: if $T_0x+V_0\eta=0$, the CNC13 injectivity
at this cutoff gives $x=0$, and the lower Vandermonde gives
$\eta=0$. Thus $\dim\mathcal S^{\rm raw}_k=q_k-v$.

For the upper endpoint set
\[
 \begin{gathered}
 P_A=\mathcal T_A\chi_j,\qquad
 (\mathcal T_Ap)(S')=\sum_{a,b=0}^8a_{ab}p(S'+\beta_{8,a,b}),\\
 R_i(S')=\mathcal T_{B_i}P_A(S'),\qquad
 (\mathcal T_{B_i}p)(S')=\sum_{a,b=0}^4b_{i,a,b}p(S'+\beta_{4,a,b}),\\
 \mathsf R_1=
   \operatorname{stack}_{i=1}^4
     \bigl(\mathcal B_k\operatorname{coeff}_{S'}R_i\bigr)^{\mathsf T},
 \qquad \mathsf R_0:\mathcal H_0\longrightarrow0 .
 \end{gathered}\tag{CIT5}
\]
All coefficient columns are padded through $D_1$.
Here $\deg P_A=q_j-v=D_1$ and
$\operatorname{lc}P_A=\mu_v\binom{q_j}{v}$. The complete LCC9--12
proof shows that $\mathsf R_1$ has rank four: the four original
$L_i^4$ are independent on the rank-four Segre quadric, and
translation of this polynomial is injective on the 25 distinct
original shifts. The native row in CIT5 is an ordinary transpose
because it evaluates the complex-linear functional on $R_i$.

The complete primitive domain is
\[
 \mathcal Y_e^{\rm raw}=\ker\mathsf R_e,\qquad
 \mathcal S_k^{\rm raw}\subset\mathcal Y_e^{\rm raw},\qquad
 r_e=\dim(\mathcal Y_e^{\rm raw}/\mathcal S_k^{\rm raw})
                  =8k+24-3e .
 \tag{CIT6}
\]
For $e=0$ it is all of $\mathcal H_0$. To prove the target
condition, at that endpoint the raw analytic image in each
$\mathcal H_1$ has dimension $q_j-v=L_0$ and fills it.
At $e=1$ its only missing coordinate is
$\Lambda_F(P_A)$ by LCC10. Therefore
$\mathsf T h$ belongs to the fourfold raw target analytic
image precisely when $\mathsf R_1h=0$. The exact numerator
cochain square kills $\mathcal S_k^{\rm raw}$ in those rows.
Subtracting $\dim\mathcal S_k^{\rm raw}=q_k-v$ from
$\dim\ker\mathsf R_e=L_e-4e$ proves CIT6.

Choose a full primitive lift matrix $P_e$ so that
$[V_0,T_0,P_e]$ is a frame of $\mathcal Y_e^{\rm raw}$.
This is an exact computable choice: first express CIT4--5 in
ordinary Taylor coordinates, take the first nonzero pivot minors
in their original order to obtain the kernel and a complement,
then transport those complement columns by $U_k^{-1}$.
It depends on the actual coefficients and imposes no metric
conditioning estimate. The chosen Taylor columns are fixed;
the original norm remains CIT2.

In the original target coefficient coordinates define
\[
 \begin{gathered}
 P_{\mathcal L_1}=\mathcal V(\mathcal V^*\mathcal V)^{-1}\mathcal V^*,
 \quad K_1=(I-P_{\mathcal L_1})\mathcal T_1,\quad
 \mathcal D_e=K_1^*K_1=\operatorname{diag}(D_{j,D_e}^{(s_1)},4),\\
 F_e=\mathcal D_e^{-1}K_1^*(I-P_{\mathcal L_1})\mathsf T P_e,
 \qquad
 B_{\rm cof}=\operatorname{stack}_{i=1}^4 M_{i,k}.
 \end{gathered}\tag{CIT7}
\]
The inverses exist by the full lower Vandermonde and CNC13.
By CIT6 the projected columns on the right belong to
$\operatorname{im}K_1$, so the displayed left inverse gives
their unique original coefficient coordinates. Consequently
\[
 K_1F_e=(I-P_{\mathcal L_1})\mathsf T P_e,\qquad
 K_1B_{\rm cof}=(I-P_{\mathcal L_1})\mathsf T T_0 .
 \tag{CIT8}
\]
These are literal original finite matrices. In particular the
$M_{i,k}$ retain the exact KRC scalar-bearing conjugation CNC12.

\subsection{Complete connecting kernel and attained target norm}

We prove the full kernel of the connecting construction.
The exact finite intersection needed here is
\[
              \operatorname{im}\mathsf T\cap\mathcal L_1
                            =\mathsf T\mathcal L_0 .
 \tag{CIT9}
\]
For a direct proof, write $\mathsf T h=\mathcal V\eta$.
Apply the six original compatibility rows. Their coefficient
square is $\mathcal R\mathcal V=\mathcal V_2\mathcal K$,
so $\mathcal V_2\mathcal K\eta=0$. At the present cutoffs
$D_e+1-q_k=8k+24-v+e\ge16$; hence the original degree-$k$
evaluation $\mathcal V_2$ is injective. The complete GKC kernel
gives $\eta=\mathcal P\xi$. The lower root-addition square says
$\mathcal V\mathcal P=\mathsf T V_0$, whence
$\mathsf T h=\mathsf T V_0\xi$. Since $\mathsf S\mathsf T=I$,
$h=V_0\xi$. The reverse inclusion is exactly that same square.
This proves CIT9 without a transversality or inverse-norm assumption.

For $h\in\mathcal Y_e^{\rm raw}$ let $c(h)$ be the coefficient
column defined as in CIT7. The native exterior differential of
$\mathsf T h$ is zero, since $B_iB_t=B_tB_i$. Replacing its
representative by lower evaluations changes its next differential
by lower evaluations, by root addition. Therefore $c(h)$ is a
cycle in the actual LCC image complex. The complete map is
\[
 \begin{gathered}
 \partial_e:\mathcal Y_e^{\rm raw}/\mathcal S_k^{\rm raw}
     \longrightarrow H^1(\mathcal I_e),\qquad
          [h]\longmapsto[c(h)],\\
 \operatorname{im}\partial_e=\mathcal C_e,\qquad
 \mathcal C_e=
 \mathscr S_e/\mathscr B,\quad
 \mathscr S_e=\operatorname{im}[B_{\rm cof},F_e],\quad
 \mathscr B=\operatorname{im}B_{\rm cof}.
 \end{gathered}\tag{CIT10}
\]
If $h=T_0x+V_0\xi$, its target coefficients are the boundary
$B_{\rm cof}x$ by CIT8. Conversely, if $c(h)=B_{\rm cof}x$,
then $\mathsf T(h-T_0x)\in\mathcal L_1$ and CIT9 gives
$h-T_0x\in\mathcal L_0$. Thus the kernel before quotient is
exactly $\mathcal S_k^{\rm raw}$. This proves injectivity and
the stated presentation; $[B_{\rm cof},F_e]$ is independent.
It is exactly the positive connecting image of
$H^0(\mathcal R_e)=\ker\mathsf t_e$ in LCC13, since its
construction lifts the original degree-zero quotient cycle
and takes its differential with the same positive sign.
No splitting of the remaining $\ker U_e$ has been introduced.

Its full cohomology restriction Gram in the $P_e$ image coordinates is
\[
 \begin{split}
 H_{{\rm conn},e}
 ={}&F_e^*\mathcal D_eF_e\\
 &-F_e^*\mathcal D_e B_{\rm cof}
       (B_{\rm cof}^*\mathcal D_eB_{\rm cof})^{-1}
                           B_{\rm cof}^*\mathcal D_eF_e .
 \end{split}\tag{CIT11}
\]
The exact least norm over every original boundary gives this
Schur complement. All original lower coefficients were already
minimized simultaneously in $\mathcal D_e$. Equivalently set
\[
 Z_e^{\rm raw}=[\mathcal V,\mathsf T T_0],\qquad
 H_{{\rm conn},e}
  =(\mathsf T P_e)^*(I-P_{Z_e^{\rm raw}})\mathsf T P_e,
 \quad P_Z=Z(Z^*Z)^{-1}Z^* .
 \tag{CIT12}
\]
The columns of $Z_e^{\rm raw}$ are independent by CIT9 and
the independence in CIT4. The same argument with the added
$P_e$ proves independence of $[Z_e^{\rm raw},\mathsf T P_e]$.
Indeed the complete minimization in CIT11 allows exactly
$\mathsf T T_0x+\mathcal V\eta$. This proves CIT12 and positivity.
It also gives the full original raw determinant
\[
 \det H_{{\rm conn},e}
 =\frac{\det([Z_e^{\rm raw},\mathsf T P_e]^*
                         [Z_e^{\rm raw},\mathsf T P_e])}
        {\det((Z_e^{\rm raw})^*Z_e^{\rm raw})}.
 \tag{CIT13}
\]
These are the same fixed combined frames, so their ordinary
Schur determinant introduces no missing coordinate factor.

\subsection{The original source-to-image determinant at its low cutoff}

Give the primitive source quotient its original source metric
and its correlated image metric, respectively:
\[
 H_{{\rm src},e}=P_e^*(I-P_{[V_0,T_0]})P_e,\qquad
 H_{{\rm cor},e}=(\mathsf T P_e)^*
                  (I-P_{\mathsf T[V_0,T_0]})\mathsf T P_e .
 \tag{CIT14}
\]
These are two explicit attained minima. The original NBG
$\|\mathsf T\|\le K_+$ and $\|\mathsf S\|\le K_-$ with
$\mathsf S\mathsf T=I$ imply
\[
 K_-^{-2}H_{{\rm src},e}\preceq H_{{\rm cor},e}
                  \preceq K_+^2H_{{\rm src},e}.
 \tag{CIT15}
\]
For each fixed primitive value, apply the two norm inequalities
to every representative in $P_eb+\mathcal S_k^{\rm raw}$ and
take the minima. This proves both directions on the exact same
fibre. The finite constants are the complete NBG12 ones at
base $k$ and cutoff $D_e$, including source order $(1,1)$ or
$(k,j)$ and center shift $(\beta_{4,a,b}-2)/i$.

There is also an exact independent-gauge quotient identity.
In the original source quotient put
\[
 X_k=\mathcal S_k^{\rm raw}/\mathcal L_0,\qquad
 Y_e=\mathcal Y_e^{\rm raw}/\mathcal L_0,\qquad X_k\subset Y_e .
\]
Their correlated target images are the original NKC9 images.
The original independent-gauge form restricted to $Y_e$ and
then quotiented by $X_k$ allows precisely the total fibre
$\mathsf T T_0x+\mathcal L_1$; its correlated quotient allows
$\mathsf T T_0x+\mathsf T\mathcal L_0$. These are CIT12 and
CIT14. Hence the literal PCJ7 source quotient gives
\[
 \frac{\det H_{{\rm conn},e}}{\det H_{{\rm cor},e}}
                =\det\Omega_{Y_e/X_k,D_e}.
 \tag{CIT16}
\]
If $R$ transports the original primitive lifts to an adapted
correlated quotient frame, both Grams are changed by $R^*(\cdot)R$,
and their two factors $|\det R|^2$ cancel in CIT16. This is
the exact coordinate map behind that determinant equality.

Use the actual compatibility-plus-complement expressions
$\chi_{X_k,D_e}$ and $\chi_{Y_e,D_e}$ of PCJ1.
The complete shared-update proof PCJ13--16 gives
\[
 \begin{gathered}
 \log\det H_{{\rm conn},e}-\log\det H_{{\rm src},e}
   =-(\chi_{X_k,D_e}-\chi_{Y_e,D_e})+\epsilon_e,\\
 -2r_e\log\max(1,K_-)-r_e\log c_+
      \le\epsilon_e
      \le2r_e\log\max(1,K_+)-r_e\log c_- .
 \end{gathered}\tag{CIT17}
\]
Indeed write the left side as the correlated/source logarithm
plus $\log\det\Omega_{Y_e/X_k}$. The first is bounded by
CIT15. PCJ16 writes the negative of the latter as
$\chi_{X_k}-\chi_{Y_e}+\nu$, with
$r_e\log c_-\le\nu\le r_e\log c_+$. Subtraction gives exactly
CIT17. Here $c_-=\min(1,L_-^{-2})$ and
$c_+=\max(1,L_+^2)$ are the original NKC constants at this
cutoff. Consequently this specified error is
$O_{\rm actual}(k\log q_k)=o(kq_k)$ in both original source pairs.
At $e=0$, $Y_0$ is the full ambient source quotient, so its
source-complement projector vanishes and $\chi_{Y_0}=0$
exactly. At $e=1$ its four actual polynomial constraints
remain; $\chi_{Y_1}$ is not set to zero.

CIT14--17 concern the defining low cutoff $D_e$. They do not
assert that a primitive truncated germ has a common compatible
extension to every later cutoff. The target coefficient columns
$F_e$ now are fixed, and their complete later metrics are
calculated next without that assertion.

\subsection{The fixed connecting flag at every original target row}

For $q_j-1\le N\le2q_j$ put $D=N-v$ and retain the fixed
$B_{\rm cof},F_e$ in the original target coefficients. Define
\[
 \begin{gathered}
 \mathcal D_N=\operatorname{diag}(D_{j,N-v}^{(s_1)},4),\qquad
 Z_{B,N}=[\mathcal V_N,\mathcal T_{1,N}B_{\rm cof}],\\
 Z_{C,N}=[Z_{B,N},\mathcal T_{1,N}F_e],\qquad
 H_{{\rm conn},e}(N)
   =\operatorname{Gram}_{\mathcal D_N}(\mathscr S_e/\mathscr B),\\
 \det H_{{\rm conn},e}(N)=\det(Z_{C,N}^*Z_{C,N})/
                                      \det(Z_{B,N}^*Z_{B,N}).
 \end{gathered}\tag{CIT18}
\]
The original DNE positivity and the independence in CIT10
make both raw frames full column rank at every one of these
indices. Eliminating the complete lower frame and then
$B_{\rm cof}$ proves the determinant identity. At the defining
$N=q_j-1+e$, it is exactly CIT11--13 by CIT8.

Every entry is given by the original arithmetic recursion.
If $\nu_{a,n}=\sum(Z_{\exp,j})_{uv,a}\beta_{j,u,v}^{\,n}$, then
\[
 \begin{gathered}
 f_{a,n}=
 \frac{\nu_{a,n+v}-\sum_{d=v+1}^{n+v}
                 \binom{n+v}{d}\mu_df_{a,n+v-d}}
      {\binom{n+v}{v}\mu_v},\\
 b_{0,m}=\mathbf1_{m=0},\quad b_{-1,m}=0,\qquad
 b_{n+1,m}=i^{-1}(b_{n,m-1}-c'_j b_{n,m})
                       -n(n-1+s_1/2)b_{n-1,m},\\
 t_n=\left(\frac{\sum_m b_{n,m}f_{a,m}}
                    {\sqrt{(2\pi)^{s_1/2}n!(s_1/2)_n}}\right)_a,
 \quad
 v_n=\left(\frac{\sum_m b_{n,m}\beta_{j-8,u,v}^{\,m}}
                    {\sqrt{(2\pi)^{s_1/2}n!(s_1/2)_n}}\right)_{u,v}.
 \end{gathered}\tag{CIT19}
\]
Arrays outside their polynomial range are zero. Differentiating
$E_AF_a=N_a$ gives the first line; substituting the full original
Gamma recurrence gives the next. Thus these are exactly the
unchanged IRR7--9 rows, retaining all original moments and phases.
The new four-row blocks in CIT18 are
$z_{B,n}=[\operatorname{diag}(v_n,4),
\operatorname{diag}(t_n,4)B_{\rm cof}]$ and
$z_{C,n}=[z_{B,n},\operatorname{diag}(t_n,4)F_e]$.
For $X=B,C$ define $G_{X,n}=Z_{X,n+v}^*Z_{X,n+v}$.
In degree-major row order each raw Gram increases by $z^*z$.
The fixed row permutation from copy-major order is unitary.

For $q_j-v\le n\le2q_j-v$, the exact connecting-volume increment is
\[
 \log\frac{\det H_{{\rm conn},e}(n+v)}
                   {\det H_{{\rm conn},e}(n+v-1)}
 =\log\det(I_4+z_{C,n}G_{C,n-1}^{-1}z_{C,n}^*)
  -\log\det(I_4+z_{B,n}G_{B,n-1}^{-1}z_{B,n}^*)\ge0 .
 \tag{CIT20}
\]
The determinant lemma proves the equality. For its sign,
transform the old extra columns to their full least-norm
section over $Z_{B,n+v-1}$. The larger row covariance is the
smaller one plus $R Q^{-1}R^*\succeq0$, by block inversion.
This proves the inequality and retains the actual old quotient
denominator. No independent rowwise gauge minimum is taken.

\subsection{Its exact place in the common low invariant loss}

Let $\mathbb J=\ker\mathbb B$ be the full original fourfold
DNE invariant kernel. At each later $N$ use the same actual
target metric $\mathcal D_N$ and fixed flags
\[
 \begin{gathered}
 W=\mathscr B+\mathbb J,\qquad
 K_{{\rm conn},e}=\mathscr S_e\cap W
                 =\mathscr B+(\mathscr S_e\cap\mathbb J),\\
 T_{{\rm conn},e}=\mathscr S_e+W=\mathscr S_e+\mathbb J,\\
 \omega_{{\rm conn},e}(N)=
 \frac{\det\operatorname{Gram}_{\mathcal D_N}
                       (T_{{\rm conn},e}/W)}
      {\det\operatorname{Gram}_{\mathcal D_N}
                       (\mathscr S_e/K_{{\rm conn},e})}.
 \end{gathered}\tag{CIT21}
\]
The two Grams use the same quotient-value coordinates under
$[z]\mapsto[\mathbb Bz]$. Its original fibres are exactly the
two spaces in the denominator and numerator, as in LCC20--23.
The source is the corresponding quotient of the connecting
cohomology image; its full kernel is
$(\mathscr S_e\cap\mathbb J)/(\mathscr B\cap\mathbb J)$.
Thus CIT21 is its actual further target loss, and is in $(0,1]$.

The full low cycles contain $\mathscr S_e$, and LCC13 gives
\[
 \dim(\mathscr Z_e/\mathscr S_e)=\kappa_e
                              \le3\max(v-e,0).
 \tag{CIT22}
\]
Both spaces are fixed original coefficient flags for the later
metric comparisons. The difference of their invariant loss
returns has the finite bound
\[
 \left|\mathcal R\log\omega_e-\mathcal R\log\omega_{{\rm conn},e}\right|
       \le2\kappa_e\log A_{j,s_1}^{\rm end}
                            =O_{\rm actual}(q_j\log j)=o(kq_k).
 \tag{CIT23}
\]
Here $\mathcal Rf=-f(q_j-1)-f(q_j)+f(2q_j-1)+f(2q_j)$.
For a complete verification, apply the actual inclusions
$\mathscr S_e\subset\mathscr Z_e$ with their common $W$.
The kernel/intersection quotient and sum quotient fit into
the exact sequence
$0\to K_{\rm full}/K_{\rm conn}\to
\mathscr Z_e/\mathscr S_e\to
T_{\rm full}/T_{\rm conn}\to0$ by the identity on representatives.
Their two ranks sum to $\kappa_e$. The full adapted-Gram
factorization LFC7--12 expresses the loss difference as the
sum of those two quotient log determinants minus the
$\mathscr Z_e/\mathscr S_e$ log determinant and a fixed
frame constant. FRE's same native form sandwich puts the
three monotone returns between zero and twice their ranks
times $\log A_{j,s_1}^{\rm end}$. Both sides of the resulting
signed difference are bounded by $2\kappa_e\log A$.
This proves CIT23, including $\kappa_e=0$, with all frame
constants canceled only after the original four signs.

There is a further exact relation to the ORIGINAL ITL source.
Write $F_{\rm inv}=\operatorname{stack}_i(M_{i,k}J_0)$.
Since $B_{\rm cof}=\operatorname{stack}_iM_{i,k}$,
\[
 \operatorname{im}F_{\rm inv}\subset\mathscr B,\qquad
 \mathbb J\subset
 W_{\rm inv}:=\mathbb J+\operatorname{im}F_{\rm inv}
          \subset W\subset T_{{\rm conn},e}.
 \tag{CIT24}
\]
These inclusions are literal original coefficient maps.
The ITL source therefore represents boundaries in low
cohomology, but its target Gram remains in the denominator
$\operatorname{Gram}(W)$ of the cohomology target quotient.

To exhibit every factor, let $I_R=F_{\rm inv}Z_R$ be exactly
the fixed ITL4 source-kernel complement and choose fixed
columns $R_b$ completing $[X_{\mathbb J},I_R]$ to $W$.
In a specified fixed full frame $X_W$, write
\[
 [X_{\mathbb J},I_R,R_b]=X_W C_W,\qquad
 Q_b(N)=\operatorname{Gram}_{\mathcal D_N}(W/W_{\rm inv})
                     \quad\hbox{in the lifts }R_b .
\]
The first quotient of this chain has exactly ITL's target
Gram $H_{\rm inv}(N)$: its fibre is $I_Rb+\mathbb J$.
Its source pullback quotient is the full ITL $A_{\rm inv}(N)$,
and $\det H_{\rm inv}=\det A_{\rm inv}\,\omega_{\rm inv}$.
Sequential Schur complements consequently give
\[
 |\det C_W|^2 g_W(N)
  =g_{\mathbb J}(N)\det H_{\rm inv}(N)\det Q_b(N),
 \qquad g_X(N)=\det(X_X^*\mathcal D_NX_X).
 \tag{CIT25}
\]
All matrices in this formula are in their stated independent
frames; $r_{\rm inv}=0$ gives the empty first quotient.

Finally retain the full LFC coordinate constant $c_{{\rm conn},e}$
from the adapted frames for
\[
 (\mathscr S_e,K_{{\rm conn},e},T_{{\rm conn},e},W).
\]
Its cross ratio and CIT25 give the exact original identity
\[
 \begin{split}
 \log\omega_{{\rm conn},e}+\log\omega_{\rm inv}
  ={}&\log g_{T_{{\rm conn},e}}+\log g_{K_{{\rm conn},e}}
       -\log g_{\mathscr S_e}-\log g_{\mathbb J}\\
    &-\log\det A_{\rm inv}-\log\det Q_b
       +\log c_{{\rm conn},e}+2\log|\det C_W| .
 \end{split}\tag{CIT26}
\]
This relates the two actual losses through their exact common
boundary structure; it does not identify their source spaces.
Taking the original four-return cancels the last two fixed
coordinate constants. Equivalently ITL16's exact matched
relative operator satisfies
$\det H_{\rm inv}=\det Q_{\rm inv}\det\mathscr R_{X}$,
where $Q_{\rm inv}$ is the full original source quotient
$J_0^*D_{k,N-v}^{(s_0)}J_0/E$. Substituting this equality
instead in CIT25 replaces
$(\omega_{\rm inv},A_{\rm inv})$ in CIT26 by
$(\det\mathscr R_X,Q_{\rm inv})$, preserving the source-order
pair $(1,1)$ or $(k,j)$. It retains the original source mass.

Thus the surviving connecting target is given by the explicit
matrix CIT7, its complete target minimum CIT11--13, the
original low source-to-image determinant CIT17, and all later
original row formulas CIT18--20. The fixed-rank remainder is
controlled by CIT23. The remaining common target loss has
the exact original ITL boundary relation CIT24--26, with its
whole residual overlap, all absolute native Grams and both
source measures still present.



% Complete original provider: ORIGINAL_PRIMITIVE_CONNECTING_TRANSPORT.tex
% Exact back-propagation of PDS into the actual CIT connecting-image metric.
\section{Primitive-source transport inside the original connecting image}
\label{sec:PDI}

Keep the complete original PDS1--24 and CIT1--26 data at fixed
$k=4l+1\ge9$, $j=k+4$, $D_e=(k+5)^2-v-1+e$, and either original
source-order pair $(s_0,s_1)=(1,1)$ or $(k,j)$. All polynomial
coefficients, period data, lower gauges and exterior signs remain
those of those constructions. In this section $P_e,F_e,T_{a,D_e},
V_{a,D_e}$ are CIT's primitive lifts, connecting target columns,
analytic frame and full lower frame. The polynomial coefficient
matrix $C_e$ and metric $G_e=\overline{(C_e^*C_e)^{-1}}$ are PDS7.
The boundary column matrix is the unchanged
$B_{\rm cof}=\operatorname{stack}_iM_{i,k}$.

\subsection{The exact source coordinate map and its target naturality}

The coordinate maps between the two proved primitive presentations are
\[
 \Theta_e=C_e^{\mathsf T}P_e,\qquad
 \Theta_0:\mathbb C^{d_0}\xrightarrow{\ \simeq\ }\mathbb C^{d_0},
 \qquad
 \Theta_1:\mathbb C^{d_0-3}\xrightarrow{\ \simeq\ }\ker W^{\mathsf T},
 \quad d_0=8k+24 .
 \tag{PDI1}
\]
The inverse signs mean inverses onto the displayed images.
Indeed PDS4 is the onto quotient by precisely CIT4's analytic
source, and $P_e$ is a full basis of that primitive quotient.
The PDS8 four-row condition is CIT5's same polynomial differential.
Thus both maps are bijections onto the asserted spaces. Equality
of the exact minimum-norm fibres proves
\[
 H_{{\rm src},e}=\Theta_e^*G_e\Theta_e .
 \tag{PDI2}
\]
For completeness a native vector $w$ and $P_ea$ have the same
PDS quotient coordinates exactly when
$C_e^{\mathsf T}(w-P_ea)=0$, equivalently
$w-P_ea\in\operatorname{im}[V_{k,D_e},T_{k,D_e}]$ by PDS4.
The minimum over this fibre is both sides of PDI2, by CIT14 and
PDS7. This proves the metric identity as well as its linear map.

Let $J_k^{\rm tr}:\mathcal H_{k,D_1}\to\mathcal H_{k,D_0}$ and
$J_j^{\rm tr}:\mathcal H_{j,D_1}\to\mathcal H_{j,D_0}$ be literal jet
truncation, so they delete the last native functional coordinate;
the source order and centre are fixed during this truncation.
Let $\operatorname{pr}_0$ delete the last polynomial quotient
coordinate as in PDS11, and set
\[
 Z=\operatorname{pr}_0\Theta_1,\qquad
 R=\Theta_0^{-1}Z,\qquad
 \operatorname{rank}R=d_0-3,\qquad
 \Theta_1=EZ .
 \tag{PDI3}
\]
The rank assertion and last equality follow from PDS11. The
first $d_0$ relation polynomials are unchanged between cutoffs;
therefore
$C_0^{\mathsf T}J_k^{\rm tr}P_1=\operatorname{pr}_0\Theta_1=\Theta_0R$.
Their difference is in the old analytic source. In its original
independent frame there are unique matrices $L$ and $A$ such that
\[
 J_k^{\rm tr}P_1=P_0R+V_{k,D_0}L+T_{k,D_0}A .
 \tag{PDI4}
\]
They are explicit finite matrices: with
$U=[V_{k,D_0},T_{k,D_0}]$ their exact block expression is
\[
 \begin{bmatrix}L\\A\end{bmatrix}
   =(U^*U)^{-1}U^*(J_k^{\rm tr}P_1-P_0R).
\]

The actual connecting coefficient columns satisfy
\[
 \boxed{\quad F_1=F_0R+B_{\rm cof}A .\quad}
 \tag{PDI5}
\]
To prove it, the original multiplier square commutes with
truncation: $(\operatorname{diag}(J_j^{\rm tr},4))\mathsf T^{(1)}
=\mathsf T^{(0)}J_k^{\rm tr}$. This is the equality of the Taylor
convolutions in CIT3 after discarding the last coefficient.
Likewise truncation takes each original lower or analytic target
column to the corresponding old column, because they are jets of
the same entire or analytic germ. Apply these equalities to the
target representative of $F_1$ in CIT8 and use PDI4. Its old
target jet differs from that of $F_0R+B_{\rm cof}A$ by an allowed
lower evaluation, using the root-addition square on $V_{k,D_0}L$.
The original degree-$j$ analytic coefficient map modulo lower
evaluations is injective at $D_0$ by CNC13. Thus the coefficient
columns themselves agree, proving PDI5. No inverse estimate is
used to pass from equality of the jets to these coefficients.

In particular, in the single original target coefficient space,
\[
 \mathscr S_1=\operatorname{im}[B_{\rm cof},F_1]
 \subset\mathscr S_0=\operatorname{im}[B_{\rm cof},F_0],\qquad
 \dim(\mathscr S_0/\mathscr S_1)=3 .
 \tag{PDI6}
\]
Independence of each displayed frame is proved in CIT10, and
PDI3 gives the rank of the quotient inclusion. This proves the
codimension and identifies the actual map from primitive
truncation to the fixed connecting coefficient flags.

\subsection{The two low metric increments in common coordinates}

The old source restricted to this image and the upper source,
in precisely the $P_1$ coordinates, have Grams
\[
 K_{{\rm src},0}=R^*H_{{\rm src},0}R=Z^*G_0Z,\qquad
 K_{{\rm src},1}=H_{{\rm src},1}=(EZ)^*G_1(EZ).
 \tag{PDI7}
\]
PDS12--24 now apply to these actual connecting-source frames:
\[
 K_{{\rm src},1}=K_{{\rm src},0}
       +\frac{(\eta Z)^*(\eta Z)}{\sigma},\qquad
 0\le\Delta_{\rm src}:=
        \log\frac{\det K_{{\rm src},1}}{\det K_{{\rm src},0}}
       \le\log U_*=o(kq_k).
 \tag{PDI8}
\]
Every constant in $U_*$ is the original PDS19 finite product;
it has the fully proved explicit bound PDS24.

In the original target define
\[
 K_{{\rm tgt},0}=R^*H_{{\rm conn},0}R,\qquad
 K_{{\rm tgt},1}=H_{{\rm conn},1} .
 \tag{PDI9}
\]
By PDI5 these are the metrics of the same fixed quotient
$\mathscr S_1/\mathscr B$ at the two native target cutoffs,
in the same $F_1$ value coordinates. Adding the boundary
$B_{\rm cof}A$ does not change its minimum over that boundary.
This proves the first identity in PDI9 as a quotient statement,
including its frame, rather than presuming equality of two
arbitrary representatives.

Here is its complete four-row update. At $D_0$ set
\[
 \begin{gathered}
 N_B=[\operatorname{diag}(V_{j,D_0},4),
             \operatorname{diag}(T_{j,D_0},4)B_{\rm cof}],\qquad
 N_F=\operatorname{diag}(T_{j,D_0},4)F_1,\\
 G_B=N_B^*N_B,\quad G_{BF}=N_B^*N_F,\quad
 K_{{\rm tgt},0}=N_F^*N_F-G_{BF}^*G_B^{-1}G_{BF} .
 \end{gathered}\tag{PDI10}
\]
Both original lower and boundary frames are kept in $N_B$.
Let $z_B,z_F$ be the next degree's four rows of these very
frames, computed from CIT19, and define
\[
 S_B=I_4+z_BG_B^{-1}z_B^*,\qquad
 r_B=z_F-z_BG_B^{-1}G_{BF} .
 \tag{PDI11}
\]
Their exact update is
\[
 \begin{gathered}
 K_{{\rm tgt},1}=K_{{\rm tgt},0}+r_B^*S_B^{-1}r_B,\\
 \Delta_{\rm tgt}:=
 \log\frac{\det K_{{\rm tgt},1}}{\det K_{{\rm tgt},0}}
 =\log\det\left(I_4+
 S_B^{-1/2}r_BK_{{\rm tgt},0}^{-1}r_B^*S_B^{-1/2}\right)\ge0 .
 \end{gathered}\tag{PDI12}
\]
For proof replace $N_F$ by its old orthogonal residual
$N_F-N_BG_B^{-1}G_{BF}$. In this unchanged quotient the old
joint Gram is $\operatorname{diag}(G_B,K_{{\rm tgt},0})$;
the added rows are $[z_B,r_B]$. The new Schur complement is
$K_{{\rm tgt},0}+r_B^*[I-z_B(G_B+z_B^*z_B)^{-1}z_B^*]r_B$.
Multiplying by $I+z_BG_B^{-1}z_B^*$ proves that the bracket
is $S_B^{-1}$. Taking the determinant gives PDI12. The column
operation has identity diagonal blocks, so its determinant is
one and it leaves no missing frame factor.

There is a finite bound using the complete original native
next-row budget. Write $\mathsf D_0=D_{j,D_0}^{(s_1)}$ for
the single-copy original coefficient metric, $T=T_{j,D_0}$,
$V=V_{j,D_0}$, and let $t,v_\mathrm{row}$ be their next
original rows. Set
\[
 \begin{gathered}
 A_V=V^*V,\quad
 \zeta=v_\mathrm{row}A_V^{-1}v_\mathrm{row}^*,\quad
 \rho=t-v_\mathrm{row}A_V^{-1}V^*T,\\
 \beta_{\rm next}=4\log\left(1+
       \frac{\rho\mathsf D_0^{-1}\rho^*}{1+\zeta}\right),\qquad
 0\le\Delta_{\rm tgt}\le\beta_{\rm next}
       \le4\log A_{j,s_1}^{\rm end}=o(kq_k).
 \end{gathered}\tag{PDI13}
\]
The exact lower-gauge Schur update is
$\mathsf D_1=\mathsf D_0+\rho^*\rho/(1+\zeta)$, as in
IRR and LFC. Thus the full four-copy ambient log-determinant
increment is $\beta_{\rm next}$. To prove the middle bound
without multiplying by a quotient dimension, put the old
ambient form into its isometric coordinates and write the
new one as $e^{B}$ with $B\succeq0$. For any fixed subspace
with full frame $Y$, differentiation along $e^{tB}$ gives
\[
 \frac{d}{dt}\log\det(Y^*e^{tB}Y)
 =\operatorname{Tr}(BP_Y(t)),\quad
 P_Y(t)=e^{tB/2}Y(Y^*e^{tB}Y)^{-1}Y^*e^{tB/2} .
\]
$P_Y(t)$ is the orthogonal projector onto $e^{tB/2}Y$.
For nested fixed spaces its two projectors are ordered, so
the derivative of their quotient log determinant is between
zero and $\operatorname{Tr}B$. Integrating proves that the
increment of $\mathscr S_1/\mathscr B$ is at most the full
ambient increment, which is $\beta_{\rm next}$. Finally
FRE's original endpoint form bound implies that each of the
four identical single-copy rank-one relative eigenvalues is
at most $A_{j,s_1}^{\rm end}$, proving the last inequality.
Its already proved growth is $\log A_{j,s_1}^{\rm end}
=O_{\rm actual}(q_j\log j)=o(kq_k)$. The isometric coordinates
in this proof establish a comparison of the unchanged forms;
the formulas PDI10--13 remain in their original frames.

The exact adjacent source-to-image determinant correction is
therefore the difference of the explicit four-row and one-row
log determinants PDI12 and PDI8:
\[
 \begin{gathered}
 \Delta_{\rm image/source}:=
 \log\frac{\det K_{{\rm tgt},1}}{\det K_{{\rm src},1}}
 -\log\frac{\det K_{{\rm tgt},0}}{\det K_{{\rm src},0}}
       =\Delta_{\rm tgt}-\Delta_{\rm src},\\
 -\log U_*\le\Delta_{\rm image/source}\le\beta_{\rm next},
 \qquad\Delta_{\rm image/source}=o(kq_k).
 \end{gathered}\tag{PDI14}
\]
The lower determinant uses the proved common primitive
subspace $R$; its three-dimensional complement has not been
silently discarded from the full old cohomology.

\subsection{The two fixed connecting invariant losses}

At every later original target cutoff keep the literal flags
of CIT21 and the same $W=\mathscr B+\mathbb J$.
By PDI6 their nesting has codimension exactly three.
The identity on representatives gives the exact sequence
\[
 0\longrightarrow K_0/K_1\longrightarrow
 \mathscr S_0/\mathscr S_1\longrightarrow
 T_0/T_1\longrightarrow0,\qquad
 K_e=\mathscr S_e\cap W,\quad T_e=\mathscr S_e+W .
 \tag{PDI15}
\]
Surjectivity follows from the definition of the sum; its
kernel consists of the class of $x\in\mathscr S_0$ with
$x\in\mathscr S_1+W$, hence of $K_0+\mathscr S_1$ modulo
$\mathscr S_1$. Its kernel presentation is $K_0/K_1$.
Thus the two end ranks sum to three.
The complete adapted Gram factorization LFC7--12 expresses
the difference of the two loss logarithms as
$\log\det\operatorname{Gram}(T_0/T_1)
+\log\det\operatorname{Gram}(K_0/K_1)
-\log\det\operatorname{Gram}(\mathscr S_0/\mathscr S_1)$
plus a fixed coordinate constant. Under the original four-return
$\mathcal Rf=-f(q_j-1)-f(q_j)+f(2q_j-1)+f(2q_j)$ that
constant cancels. FRE bounds each monotone quotient return
between zero and twice its rank times
$\log A_{j,s_1}^{\rm end}$. The positive two ranks sum to
three and the negative rank is three; therefore
\[
 \left|\mathcal R\log\omega_{{\rm conn},0}
        -\mathcal R\log\omega_{{\rm conn},1}\right|
 \le6\log A_{j,s_1}^{\rm end}=o(kq_k).
 \tag{PDI16}
\]
This is the complete primitive-image refinement of the
full-cohomology comparison in LFC. The retained
$\ker U_e$ constituents of the full LCC sequence are still
present, with their explicit CIT23 bounds. The common long
connecting loss continues to be the exact target expression
CIT21 and its ITL boundary factorization CIT24--26. The
original low source transport now feeds that target calculation
through the explicit coefficient identity PDI5, the two
native increments PDI12--14, and the proved fixed-flag
comparison PDI16.

\end{document}
